% Book — plantilla de prosa larga para texforge (formato de 6x9 pulgadas).
% Tectonic es XeTeX: la fuente ya es UTF-8, sin paquete inputenc.
% Reemplace cada texto de ejemplo y cada comentario con % por su propio contenido.
\documentclass[11pt,openany]{book}

% --- Preámbulo (solo paquetes del bundle de Tectonic; sin inputenc/listings/minted) ---
% % El idioma sale de {{language}} (por defecto spanish).
\usepackage[{{language}}]{babel}
\usepackage[T1]{fontenc}
\usepackage{lmodern}
\usepackage{microtype}
\usepackage{csquotes}
\usepackage[paperwidth=6in,paperheight=9in,inner=2.2cm,outer=1.7cm,top=2cm,bottom=2cm,headheight=14pt,headsep=12pt,footskip=14pt]{geometry}
\usepackage{fancyhdr}
\usepackage{epigraph}
\usepackage{titlesec}
\usepackage{booktabs}
\usepackage{graphicx}
\usepackage{float}
\usepackage{xcolor}
\usepackage{hyperref}

% --- Enlaces en negro: el libro se imprime en blanco y negro ---
\hypersetup{colorlinks=true, linkcolor=black, urlcolor=black, citecolor=black,
  breaklinks=true, pdftitle={{{title}}}, pdfauthor={{{author}}}}

% --- Epígrafes alineados a la derecha, discretos y sin regla ---
\setlength{\epigraphwidth}{0.60\textwidth}
\renewcommand{\epigraphrule}{0pt}

% --- Títulos de capítulo compactos para cuidar el número de páginas ---
\titleformat{\chapter}[display]
{\normalfont\huge\bfseries}{\chaptertitlename\ \thechapter}{16pt}{\Huge}
\titlespacing*{\chapter}{0pt}{-12pt}{16pt}

% --- Tipografía cómoda: párrafos con sangría, aire desde el interlineado ---
\linespread{1.06}
\setlength{\parindent}{1.5em}
\setlength{\parskip}{0pt}

% --- Cabeceras: capítulo en pares, sección en impares, folio exterior ---
% % Mantenga {{title}} corto para que la cabecera no se desborde.
\pagestyle{fancy}
\fancyhf{}
\fancyhead[LE,RO]{\thepage}
\fancyhead[RE]{\nouppercase{\leftmark}}
\fancyhead[LO]{\nouppercase{\rightmark}}
\renewcommand{\headrulewidth}{0.4pt}
% % Las páginas de apertura de capítulo conservan el folio exterior, sin regla.
\fancypagestyle{plain}{%
  \fancyhf{}%
  \fancyhead[LE,RO]{\thepage}%
  \renewcommand{\headrulewidth}{0pt}%
}

\begin{document}

  \frontmatter

  % --- Anteportada: solo el título, una página ---
  \begin{titlepage}
    \thispagestyle{empty}
    \vspace*{\fill}
    \centering
    % % La anteportada queda sobria: no agregue nada más aquí.
    {\Large {{title}} \par}
    \vspace*{\fill}
  \end{titlepage}

  % --- Portada: título, subtítulo, autor, editorial y año ---
  \begin{titlepage}
    \centering
    % % Reemplace por el título del libro.
    {\Huge\bfseries {{title}} \par}
    \vspace{0.5cm}
    % % Reemplace por el subtítulo del libro.
    {\Large {{subtitle}} \par}
    \vspace{0.8cm}
    \rule{0.3\textwidth}{0.4pt}
    \vspace{0.8cm}
    % % Reemplace por el nombre del autor o la autora.
    {\Large {{author}} \par}
    \vfill
    % % Reemplace por su editorial y el año de edición.
    {{publisher}} \par
    {{year}} \par
  \end{titlepage}

  % --- Página de créditos: una página, sin número visible ---
  \thispagestyle{empty}
  % % Reemplace este bloque con los datos de su edición.
  \noindent © {{year}} {{author}}. \par
  \noindent Primera edición: {{publisher}}, {{year}}. \par
  \vspace{0.4cm}
  \noindent Todos los derechos reservados. Ninguna parte de este libro puede
  reproducirse sin el permiso escrito del titular de los derechos. \par
  \clearpage

  % --- Dedicatoria: una línea genérica alineada a la derecha ---
  \thispagestyle{empty}
  \vspace*{\fill}
  \begin{flushright}
    % % Reemplace por su dedicatoria en una o dos líneas.
    \textit{Para quien lea de noche y subraye en silencio.}
  \end{flushright}
  \vspace*{\fill}
  \clearpage

  % --- Índice general en su propia página ---
  \tableofcontents
  \clearpage

  % --- Prefacio: presente el libro en dos párrafos breves ---
  \chapter*{Prefacio}
  \addcontentsline{toc}{chapter}{Prefacio}
  \markboth{Prefacio}{Prefacio}
  % % Reemplace por el origen de su libro: por qué existe y para quién es.
  Este libro reúne tres capítulos breves sobre la forma en que miramos los
  días comunes. Nace de notas tomadas con calma y quiere leerse sin prisa,
  una sección cada vez, con la atención puesta en lo concreto.
  % % Reemplace por la invitación al lector: cómo leer y qué esperar.
  Cada capítulo abre con un epígrafe y cierra con una pausa. Si alguna frase
  le resulta cercana, anótela al margen: este libro espera esa conversación.

  \mainmatter

  % --- Capítulos: cada uno abre en página nueva ---
  % % Reemplace por el título de su capítulo.
\chapter{El camino}

% % Reemplace por un epígrafe genérico y su atribución.
\epigraph{Caminar despacio también es avanzar.}{Atribuido a un refrán popular}

% % Reemplace por la escena inicial: dónde empieza el capítulo.
La mañana empieza con una taza caliente sobre la mesa y una lista corta de
tareas que cabe en una tarjeta. El primero en levantarse abre la ventana,
mira la calle vacía y decide que hoy basta con hacer bien tres cosas. Como
dice la casa, \enquote{lo poco, hecho con cuidado, rinde el doble}.

% % Reemplace por el desarrollo: qué cambia durante el día.
A media mañana el plan se tuerce un poco: una visita breve, una llamada
larga y un recado que nadie esperaba. En lugar de correr, el cuaderno
recibe una línea nueva y la tarde se ordena otra vez. El truco está en
mover solo lo necesario y dejar lo demás para mañana sin culpa.

\begin{center}
  * * *
\end{center}

% % Reemplace por el cierre: qué queda aprendido al final del día.
Por la noche la lista muestra dos marcas hechas y una pendiente. No es una
victoria grande, pero es una medida honesta del día. La última línea del
cuaderno dice que mañana el camino sigue, con los mismos zapatos y con un
poco más de calma.

  % % Reemplace por el título de su capítulo.
\chapter{La pausa}

% % Reemplace por un epígrafe genérico y su atribución.
\epigraph{Detenerse a tiempo ordena la marcha.}{Atribuido a un dicho común}

% % Reemplace por la escena inicial: qué se detiene y por qué.
A mitad de la semana el trabajo pide una pausa y el cuerpo la agradece. Se
cierra el cuaderno, se apaga la lámpara y se sale a caminar sin destino
fijo. El barrio enseña su ritmo lento: una tienda abierta, un perro
dormido y dos vecinos que hablan sin prisa. Entonces queda claro que,
\enquote{parar también es parte del plan}.

% % Reemplace por el método de su pausa: los pasos que sigue.
\section{ Tres pasos para ordenar la pausa}

% % Reemplace por su propio flujo; mantenga el estilo y el título.
\begin{mermaid}[style=monochrome, width=0.85\linewidth, caption={De la idea a la forma final}, pos=H]
  flowchart LR
  A[Idea] --> B[Borrador] --> C[Final]
\end{mermaid}

% % Reemplace por lo que anota en su cuaderno durante la pausa.
La pausa sigue tres pasos simples: anotar lo pendiente, elegir lo primero
y soltar el resto. El diagrama resume ese orden sin adornos.

\begin{center}
  * * *
\end{center}

% % Reemplace por su fragmento breve: un poema o una nota de registro.
% % Reemplace por su fragmento; máximo doce líneas cortas.
\begin{code}[lang=text]
nota del martes, sin prisa:
la mesa queda libre,
la taza vuelve a su lugar,
la lista espera en calma,
la tarde sigue su curso,
mañana será otro día
\end{code}

% % Reemplace por el cierre: qué trae de vuelta la pausa.
Al volver, la mesa está igual pero la cabeza está distinta. Lo urgente
sigue ahí, mas ya no manda. La pausa devuelve el orden y el trabajo
continúa con mejor pulso.

  % % Reemplace por el título de su capítulo.
\chapter{El regreso}

% % Reemplace por un epígrafe genérico y su atribución.
\epigraph{Volver a casa es empezar otra vez.}{Atribuido a una copla anónima}

% % Reemplace por la escena inicial: el camino de vuelta.
El regreso empieza con la maleta ligera y el billete en el bolsillo. Por
la ventana pasan campos cortos y casas bajas que parecen contar los
minutos. Nadie espera nada grande: solo llegar, saludar y sentarse a la
mesa. Y en ese gesto simple cabe la frase que sostiene la casa,
\enquote{lo bueno tarda, pero llega}.

% % Reemplace por el desarrollo: qué encuentra al volver.
Al llegar, todo está en su lugar aunque algo ha cambiado de sitio. La
cocina huele a pan, el patio guarda la misma silla y el reloj marca la
hora de siempre. Las noticias se cuentan despacio, sin orden, mientras la
tarde cae sobre los tejados.

\begin{center}
  * * *
\end{center}

% % Reemplace por el cierre: qué se queda después del regreso.
Por la noche no hay discurso ni promesa, solo un acuerdo claro: cuidar lo
cercano y volver a empezar cada día. El libro se cierra aquí, con la casa
en silencio y la lámpara encendida para quien llegue después.


  \backmatter

  % --- Agradecimientos: cierre breve del libro ---
  \chapter{Agradecimientos}
  % % Reemplace por sus agradecimientos en dos párrafos breves.
  Este libro debe mucho a las personas que leyeron los borradores y dijeron
  la verdad con cariño. Gracias a quienes prestaron tiempo, silencio y
  preguntas buenas cuando el texto aún no sabía lo que quería decir.
  % % Reemplace por el cierre: qué queda abierto y qué sigue.
  Gracias también a quienes cuidan las casas, las mesas y las mañanas donde
  se escribió. Lo que sigue ahora pertenece al lector: que estas páginas
  acompañen y dejen una marca amable.

\end{document}
